\subsection{The $f_3$-stress wobble-breaking translation readout}
\label{sec:q6-f3-stress-wobble-broken-translation-readout-cross-layer}

\paragraph{Finite coordinate.}
This record isolates a single BioReality conjecture: the $B^*_{Q6}$ $f_3$-stress residual coordinate is allowed to enter biological interpretation only through a matched translation-layer row. The carrier is the finite CDS-by-organism residual coordinate surface restricted to the $f_3$-stress coordinate and to translation readouts that remain paired with the same organism, CDS identifiers, and control covariates. The coordinate is not a direct claim about protein output or phenotype. It is a finite row relation among synonymous codon imbalance after amino-acid projection, organism-scoped tRNA or wobble weights, and controlled translation readouts such as tAI, stAI, ribosome occupancy, or translation efficiency.

\paragraph{Saturation datum.}
The available $f_3$ saturation row fixes the codon boundary used by this conjecture. The mediated set is
$M = \{\mathrm{AAA},\mathrm{AAG},\mathrm{ACA},\mathrm{ACG},\mathrm{AGA},\mathrm{AGG},\mathrm{AUA},\mathrm{AUG},\mathrm{CUA},\mathrm{CUC},\mathrm{CUG},\mathrm{CUU},\mathrm{UAA},\mathrm{UAG},\mathrm{UCA},\mathrm{UCG},\mathrm{UGA},\mathrm{UGG},\mathrm{UUA},\mathrm{UUG}\}$.
The saturated $f_3$ image of the $R$ set is the same list with $\mathrm{ACA}$ and $\mathrm{ACG}$ removed:
$\{\mathrm{AAA},\mathrm{AAG},\mathrm{AGA},\mathrm{AGG},\mathrm{AUA},\mathrm{AUG},\mathrm{CUA},\mathrm{CUC},\mathrm{CUG},\mathrm{CUU},\mathrm{UAA},\mathrm{UAG},\mathrm{UCA},\mathrm{UCG},\mathrm{UGA},\mathrm{UGG},\mathrm{UUA},\mathrm{UUG}\}$.
The reported $f_3$ boundary is empty. Thus the coordinate is not being promoted from a residual unmatched edge; the finite row begins from a saturated third-position contrast in which the absent boundary is itself part of the constraint.

\paragraph{Wobble-breaking condition.}
The conjectural trigger is local asymmetry of the translation weight. For a synonymous codon $c$ and its third-position paired codon $c^{f_3}$, the translation row is permitted to become non-vacuous only in contexts where $w_c \neq w_{c^{f_3}}$. If $w_c = w_{c^{f_3}}$ across the relevant codon family, the $f_3$ coordinate remains a codon-topology distinction and has no translation-dynamical content in this chapter. The survival entry is therefore written as a controlled residual score rather than as an unconditional effect:
$$
\begin{aligned}
S_{f3,T_j} > 0
\end{aligned}
$$
only after amino-acid composition, organism, GC3, length, $M$-density, and mRNA controls have been applied. A rank-degenerate all-$1.0$ matrix is a null boundary for this purpose, because it contains no surviving contrast after the controls.

\paragraph{Existing powered rows.}
Several finite rows constrain the possible interpretation before the mandatory translation score is observed. The per-organism axis-strength experiment computed $11$ organisms out of $16$ requested, with $5$ folds, at least $500$ proteins per organism, and $41$ synonymous design columns. In that panel, $d_{\mathrm{resid4}}$ led the captured spectra on average, $d_{\mathrm{tRNA}}$ followed, and the per-organism dominant count was $\{d_{\mathrm{resid4}}:11\}$. The domain counts were imbalanced, with $1$ archaeal, $2$ bacterial, and $8$ eukaryotic organisms, so domain means are descriptive rather than a balanced clade conclusion. The largest GC3 rank effect reported in that row was $d_{\mathrm{usage}}$ with $\rho = 0.136$.

\paragraph{Mechanism rows already in force.}
The tRNA-poverty boundary row tested $40$ organisms out of $46$ requested, skipped $6$, used $1000$ permutations at $\alpha = 0.05$, and called the $R$ sense codons systematically tRNA-poor under GtRNAdb tGCN plus dos Reis wobble weights. This supports a tRNA-supply explanation for $R$ avoidance, but it is still not a measured translation-efficiency survival row for the present conjecture. The universal-core tRNA-adaptation row, the non-tRNA residual optimality row, and the two-axis closure row each used $11$ computed organisms out of $16$ requested, $5$ folds, at least $500$ proteins per organism, and $41$ synonymous design columns. The non-tRNA residual optimality row identified $f_3$ translation selection as the dominant residual axis. The two-axis closure row concluded that a third axis was present because the orthogonal direction had systematically positive held-out $R^2$ and same-signed organism directions.

\paragraph{Third and fourth axes.}
The third-axis identity row identified $d_{\perp}$ with usage frequency within family, with mean cosine $0.128$ and one-sided sign-test $p = 0.000488$, again on $11$ computed organisms out of $16$ requested, $5$ folds, at least $500$ proteins per organism, and $41$ synonymous design columns. The mechanism row reported selection rather than GC as the surviving explanation after GC orthogonalization, but also stated that the CSC or positional identity still needs a separate contact. The three-axis sufficiency row concluded that a fourth mechanism is present. The fourth-mechanism probes then ruled out codon-pair identity in the codon-pair assay, excluded mRNA stability as the dominant fourth-axis identity in the yeast half-life assay, and found only a partial elongation-dwell signal in \emph{Saccharomyces cerevisiae}. These rows justify searching for a translation readout, but none of them by itself supplies the required $S_{f3,T_j}$ survival entry.

\paragraph{Readout protocol.}
For each organism and each matched CDS row, the admissible readout is a partial contribution of the $f_3$ coordinate to a translation-layer measurement after the baseline covariates have already accounted for amino-acid composition, organism identity, GC3, length, $M$-density, and mRNA abundance. The practical translation contacts are organism-scoped codon usage, tRNA-Leu gene-copy or wobble-weight summaries, tAI or stAI weights, ribosome-profiling translation efficiency, and matched mRNA abundance for computing TE as RPF divided by mRNA. A modelled tAI or stAI row is bounded evidence about codon adaptation weights; it is not an observed ribosome-footprint measurement unless the ribosome-profiling contact is separately present.

\paragraph{Falsification rule.}
The conjecture is supported only if at least one controlled translation entry is strictly positive after simpler-codon-baseline ablation. It is falsified at the translation interface if all controlled entries are non-positive or if the available readout matrix is rank-degenerate across a powered organism panel. In symbols, the null condition is
$$
\begin{aligned}
\forall j,\quad S_{f3,T_j} \leq 0,
\end{aligned}
$$
with the quantifier ranging only over matched translation readouts that pass the identifier, organism, and covariate controls.

\paragraph{Local interpretation boundary.}
This chapter does not infer mature protein abundance, folding, physical admissibility, localization, molecular function, phenotype, mechanism universality, or a general biological law from the residual coordinate. It also does not infer translation realization from codon topology alone, from an all-$1.0$ rank-degenerate matrix, or from a tAI or stAI model without the required controlled survival row. Any statement about translation dynamics needs the translation contact; any statement about folding, abundance, function, phenotype, or universality needs an additional contact at that layer.
